\documentclass[a4paper]{article}
\usepackage{MMVII}

\newcommand{\TheCmd}{OriPoseEstimRelAllPairs}

\begin{document}

\CmdTitle

\tableofcontents


 %=================================================================================

\SecDescr

This command compute the relative orientation of all 
pairs selected by \FileSty{OriPoseSelecAllPAir} using 
\FileSty{OriPoseEstimRel2Im}. As other command relative
to relative orientation, it will generally be called
impliciely by \FileSty{OriPoseEstimRelGlob}. But it can
be called explicitely for didactic reason, tuning oe fine
parametrizatioN.

 %=================================================================================

\section{Parametrisation}

The parametrisation is almost the same than \FileSty{OriPoseEstimRel2Im},
see this command.
The main difference is that :

\begin{itemize}
  \item instead of the name of $2$ images, it take the localisation
       of the result of \FileSty{OriPoseSelecAllPAir}, to run relative orientation
        on all selected pairs, it will be used also for saving the result;

  \item it does not take an output orientation \FileSty{[Ori,Out]}.
\end{itemize}


%=================================================================================

\SecRelatedCmd

\begin{itemize}
   \item  \FileSty{OriPoseEstimRel2Im}
   \item  \FileSty{OriPoseSelecAllPAir}
   \item  \FileSty{OriPoseEstimRelGlob}
\end{itemize}



\end{document}






